\section{Detailed Experimental Settings} 
\label{app:details}  
\subsection{Main Experiment Settings}
\paragraph{Models \& Dataset.} All experiments are conducted on five backbone language models: Llama-3.2-3B-Instruct~\citep{grattafiori2024llama3herdmodels}, Qwen2.5-Math-7B~\citep{yang2024qwen25mathtechnicalreportmathematical}, Qwen3-Base-1.7B and 8B~\citep{yang2025qwen3technicalreport}, and DeepSeek-7B~\citep{shao2024deepseekmathpushinglimitsmathematical}. DeepSeek-7B denotes the Qwen2.5-Math-7B model distilled from DeepSeek-R1. Following model specifications, we use a context length of 4k tokens for Qwen2.5-Math-7B, which is the maximum supported by this model, and 8k tokens for all other models. All models are trained on the DeepScaleR mathematics dataset~\citep{sun2025improvingdataefficiencyllm}. For each training instance, we use a unified instruction-style prompt template. The model is instructed to reason step by step and return the final answer enclosed in \verb|\boxed{}| tags, which enables reliable answer extraction during training and evaluation.  
\begin{tcolorbox}[ colback=gray!5, colframe=black, coltitle=white, title=\textbf{Prompt Template}, fonttitle=\bfseries, arc=6pt, boxrule=1pt, left=8pt, right=8pt, top=6pt, bottom=6pt ] Please solve the following math problem: \{\{Question Description\}\}. The assistant first thinks about the reasoning process step by step and then provides the user with the answer. Return the final answer in \verb|\boxed{}| tags, for example \verb|\boxed{1}|. Let's solve this step by step. 
\end{tcolorbox}  
\vspace{-0.6em}
\paragraph{Training.} All methods are implemented in verl~\citep{Sheng_2025}, with vLLM~\citep{kwon2023efficientmemorymanagementlarge} for rollout generation and FSDP~\citep{zhao2023pytorchfsdpexperiencesscaling} for actor optimization. Policy optimization is performed on the actor network only. We use AdamW with a constant learning rate of $1\times10^{-6}$, $\beta_1=0.9$, $\beta_2=0.999$, and weight decay 0.01. Following prior work~\citep{yu2025dapoopensourcellmreinforcement, zheng2025prosperitycollapsefaroffpolicy}, we omit the explicit KL-divergence regularization term. Rewards are binary: 1.0 if the extracted answer is correct, and 0.0 otherwise.  For rollout generation, we use temperature 1.0 and sample 8 responses per prompt. Each configuration is trained for 4096 model updates with batch size 32, giving the same rollout-sample budget across methods. We evaluate GRPO under both low- and high-staleness settings, $\mu=4$ and $\mu=1024$, and set $\mu=1024$ for \ourmethod{} and M2PO. The high-staleness GRPO baseline uses the same two-phase large-stage rollout pipeline as \ourmethod{}, but optimizes the standard GRPO objective with default clipping and no negative-advantage veto. This isolates the effect of the proposed stabilization mechanism from staged rollout execution. Thus, \ourmethod{} uses four rollout stages, each followed by 1024 model updates.  For GRPO, we set $\epsilon=0.2$. For \ourmethod{}, we use the sequence-level negative-advantage veto scope $\mathcal S_{\mathrm{seq}}$, a conservative superset of the harmful suffix $\mathcal H_\kappa$ identified in Section~\ref{sec:diagnosis}, and set $\tau_c=10^{-4}$; Appendix~\ref{app:nav_scope_ablation} ablates alternative scopes. For M2PO, we use the recommended second-moment threshold of 0.04.  

\paragraph{Efficiency Measurement.} All efficiency experiments use identical $4\times$H200 GPU configurations without parameter offloading. We report two timing metrics. $T_{\mathrm{rollout}}$ denotes the cumulative wall-clock time of all rollout-generation phases across training. $T_{\mathrm{total}}$ denotes the end-to-end wall-clock training time, including rollout generation, policy optimization, synchronization, data movement, and scheduling overhead. This aligned timing protocol is used for all methods in \cref{tab:performance_comparison}. 

\paragraph{Evaluation.} We evaluate trained models on AMC23/24~\citep{aops_amc}, AIME24/25~\citep{aops_aime}, and Math500~\citep{hendrycks2021measuringmathematicalproblemsolving, lightman2023letsverifystepstep}. Models are evaluated every 200 model updates, and we report the checkpoint with the best average performance across all benchmarks. Evaluation uses pass@1 accuracy. For AIME24/25, pass@1 is averaged over 16 sampled generations; for AMC23/24 and Math500, pass@1 is averaged over 4 sampled generations. All evaluations use the same context length as training. 

\subsection{Additional Experiment Settings}
\paragraph{Coding Experiment Details.}
For the coding experiment, we train Qwen2.5-Coder-7B~\citep{hui2024qwen25codertechnicalreport}
on the \texttt{code\_contests} dataset~\citep{sun2026improvingdataefficiencyllm}
and evaluate on LiveCodeBench~\citep{jain2024livecodebenchholisticcontaminationfree}.
We use an 8k context length and report pass@4 accuracy. The training setup follows
the main experiments except that the per-update batch size is 16. Both GRPO and
\ourmethod{} are trained for 4096 updates; \ourmethod{} uses four rollout stages
with $\mu=1024$. Wall-clock time is measured with the aligned protocol on
$4\times$H200 GPUs.

\paragraph{Asynchronous Experiment Details.}
For the asynchronous experiment, we use the fully asynchronous policy trainer in verl, where rollout workers and trainers are placed on separate devices and run concurrently. We use $4\times$H200 GPUs, with two allocated to rollout workers and two to trainers. Rollout samples are streamed through a message queue, and policy parameters are periodically synchronized from the trainer to the rollout workers. We synchronize parameters every 4 trainer updates for both async runs. Async GRPO uses no additional rollout staleness allowance, while async $\mu$-GRPO allows up to 255 synchronization intervals of stale rollouts. This gives a maximum rollout lag of $(1+255)\times 4=1024$ updates, matching the $\mu=1024$ synchronous setting. Trainer idle ratio is measured as the fraction of trainer step time spent waiting for and assembling rollout samples.